\chapter{Introduction}
\label{ch:introduction}
\addcontentsline{toc}{chapter}{Introduction}

\acrfull{y90} \acrfull{sirt}, also referred to as \acrfull{tare}, is an established treatment for primary and secondary liver malignancies \cite{sundram_selective_2017}. In this procedure, \gls{y90}-labelled microspheres are administered into the hepatic artery to preferentially irradiate tumour tissue while sparing healthy liver parenchyma. The administered activities are high (typically of the order \SI{2}{\giga\becquerel}), and clinically relevant absorbed-dose heterogeneity can arise at multiple spatial scales due to vascular delivery, microsphere clustering, and patient-specific hepatic anatomy \cite{strigari_evidence_2014}. Tumour doses of clinical interest often exceed \SI{200}{\gray} \cite{salem_clinical_2019}, and recent evidence has shown potential benefit from additional subsequent treatment, which could result in more dose to healthy liver tissue without proper treatment planning and verification \cite{reincke_hepatic_2022}. As a consequence, post-therapy quantitative imaging and dosimetry are central to treatment verification, outcome modelling, and the optimisation of any subsequent therapy.

Despite its importance, routine post-\gls{sirt} dosimetry is not universally performed in clinical practice, primarily due to practical workflow constraints and limitations in quantitative accuracy \cite{kappadath_quantitative_2024, gear_eanm_2023}. Post-therapy imaging of \gls{y90} is inherently challenging: image quality is often degraded by low effective count statistics and by physics that is less favourable than standard diagnostic radionuclide imaging, leading to increased noise and modelling error \cite{kappadath_quantitative_2024}. More broadly, emission tomography is vulnerable to substantial degradation from Poisson counting statistics as well as inaccuracies in system modelling and hardware limitations \cite{bai_tumor_2013}. Improving reconstruction quality in this low-count, model-mismatch regime is therefore a key enabling step for accurate absorbed-dose estimation and for robust patient-specific outcome prediction \cite{davis_personalisation_2020}.

In post-\gls{sirt} imaging, the two clinically relevant emission modalities are complementary. \gls{y90} \acrfull{pet} offers comparatively high spatial resolution, but is severely count-limited because of the extremely low internal pair-production branch. Bremsstrahlung \acrfull{spect}, by contrast, provides higher count statistics but suffers from poor spatial resolution and more severe modelling challenges, particularly for scatter, septal penetration, and additive corrections \cite{kappadath_quantitative_2024}. In contemporary workflows, these emission images are routinely paired with \acrfull{ct}, which provides attenuation correction and anatomical context. This naturally motivates reconstruction approaches in which information is shared across modalities, with \glsentryshort{ct} contributing structural guidance, and \glsentryshort{pet} and \glsentryshort{spect} contributing complementary functional information. Synergistic methods, where such information is incorporated during reconstruction rather than only after image formation, have gained prominence in related hybrid imaging contexts \cite{arridge_overview_2021}. Within post-\gls{y90} \gls{sirt} imaging, these approaches have the potential to improve lesion-level quantification and absorbed-dose estimation without increasing acquisition time or administered activity.

The central aim of this thesis is therefore to develop and evaluate a synergistic reconstruction framework tailored to post-\gls{y90} \gls{sirt}, with a particular focus on improving the reliability of quantitative emission imaging and therefore downstream dosimetry under clinically realistic count levels and modelling limitations. To this end, the thesis develops practical improvements to the bremsstrahlung \gls{spect} system model, introduces a weighted directional total nuclear variation prior for joint \glsentryshort{pet}/\glsentryshort{spect} reconstruction with \glsentryshort{ct}-derived directional guidance, and investigates optimisation and preconditioning strategies for this coupled penalised-likelihood problem.

The thesis is organised as follows:
\begin{itemize}
    \item \textbf{Chapter 1} introduces the \gls{pet}/\gls{spect}/\gls{ct} imaging problem, outlining the fundamental principles of each modality before describing the specific challenges of post-\gls{y90} imaging.
    \item \textbf{Chapter 2} describes the inverse problem in emission tomographic reconstruction and reviews regularisation strategies and efficient optimisation methods.
    \item \textbf{Chapter 3} addresses the specific challenges of post-\gls{sirt} bremsstrahlung imaging and develops practical improvements to the reconstruction system model.
    \item \textbf{Chapter 4} studies anatomy-informed regularisation and optimisation in a controlled post-reconstruction deconvolution setting.
    \item \textbf{Chapter 5} introduces the weighted directional total nuclear variation prior and its differentiable formulation.
    \item \textbf{Chapter 6} develops the optimisation machinery used with \acrfull{wdtnv}, including prior-aware preconditioning, subset selection, and multi-bed \glsentryshort{pet} handling.
    \item \textbf{Chapter 7} evaluates \glsentryshort{wdtnv} for joint \gls{y90} \glsentryshort{pet}/\glsentryshort{spect} reconstruction, compares it with the current state-of-the-art method \acrfull{shkem}, and discusses the main limitations and future directions.
    \item \textbf{The Appendix} documents additional software and reproducibility material that supports the thesis workflow.
\end{itemize}
